\section{Numerical implementation and use}
\label{sec:numerical-implementation}

\subsection{Code organization}

The implementation is divided into small files with one primary
responsibility each.

\begin{table}[H]
\centering
\caption{Numerical RG source files.}
\label{tab:source-files}
\begin{tabularx}{\textwidth}{@{}>{\ttfamily}l X@{}}
\toprule
\textrm{File} & \textrm{Responsibility}\\
\midrule
numerical\_rg.jl
  & Loads the constants, running coupling, splitting functions, and numerical
    solver modules. This is the entry file for users.\\
grids.jl
  & Builds the paired \(\ell_i\), \(b_i\), \(b_{*,i}\), \(\mu_i\), and \(t_i\)
    arrays in the \texttt{LatticeGrid} structure, and samples boundary data.\\
components.jl
  & Builds \texttt{SplittingKernelTable} objects at node or midpoint times,
    including regular, delta, and plus-distribution contributions.\\
convolution.jl
  & Evaluates the quark and gluon right-hand sides from precomputed kernel
    rows and jet vectors. Each convolution is a serial SIMD dot product.\\
solver.jl
  & Performs activation, upper-region evolution, lower-region
    descending-column evolution, and constructs the solution object.\\
interpolation.jl
  & Defines \texttt{StepwiseRGSolution}, domain checks, boundary-scale
    interpolation, and triangle-aware interpolation of both channels.\\
\bottomrule
\end{tabularx}
\end{table}

\subsection{Solver workflow}

A call to \texttt{solve\_jet\_rg} performs the following operations:
\begin{enumerate}[leftmargin=1.5em]
  \item Build the descending uniform log-\(b\) lattice and derive one time node
  from each value of \(b_*(b)\).
  \item Evaluate the user boundary function once at every lattice node.
  \item Evaluate \(\alpha_s(\mu_n)\) and precompute the full splitting table
  for every time and shift. For RK2, also build the midpoint table.
  \item March upward in time, activating each node on its boundary point.
  \item March downward by columns to fill the lower triangle.
  \item Store both \(N\times N\) channel histories and return a callable
  \texttt{StepwiseRGSolution}.
\end{enumerate}

\subsection{Minimal use}

The following example matches the current public API. The user-defined
\texttt{b\_sat} controls the saturation of \(b_*(b)\); the separate repository
variable \texttt{b0} remains the QCD constant
\(\bzero=2\ee^{-\gamma_E}\) used internally in \(\mu_*=\bzero/b_*\).

\begin{lstlisting}
include("Numerical RG/numerical_rg.jl")

b_sat = 1.0 # GeV^-1

bstar_func(b) = b / sqrt(1.0 + b^2 / b_sat^2)

boundary_func(; b, bstar, mu_start) = (
    exp(-2.0 * b),        # Jq(b, mu_start)
    0.5 * exp(-3.0 * b), # Jg(b, mu_start)
)

solution = solve_jet_rg(
    n_nodes = 1000,
    b_min = 0.001,
    b_max = 30.0,
    bstar_func = bstar_func,
    boundary_func = boundary_func,
    order = 2,
    method = :rk2,
)
\end{lstlisting}

The boundary function is intentionally generic. It receives the physical
arguments \texttt{b}, \texttt{bstar}, and \texttt{mu\_start}, and must return
the finite pair \texttt{(jq, jg)}. A perturbative boundary calculation and a
nonperturbative model can be inserted later without changing the solver.

\subsection{Inputs}

\begin{table}[H]
\centering
\caption{Public keyword arguments of \texttt{solve\_jet\_rg}.}
\label{tab:solver-inputs}
\begin{tabularx}{\textwidth}{@{}>{\ttfamily}l >{\raggedright\arraybackslash}p{0.18\textwidth} X@{}}
\toprule
\textrm{Argument} & \textrm{Allowed value} & \textrm{Meaning}\\
\midrule
n\_nodes
  & integer \(\ge2\)
  & Number \(N\) of paired \(\ell\) and \(t\) nodes.\\
b\_min
  & positive Float64
  & Smallest physical \(b\), in \(\mathrm{GeV}^{-1}\), represented by the
    lattice.\\
b\_max
  & Float64 \(>\texttt{b\_min}\)
  & Largest represented \(b\), in \(\mathrm{GeV}^{-1}\), beyond which zero
    closure is imposed.\\
bstar\_func
  & function
  & Maps each \(b_i\) to a finite positive \(b_{*,i}\). It must produce
    strictly increasing \(t_i=2\ln(\bzero/b_{*,i})\) in lattice order.\\
boundary\_func
  & function
  & Returns \texttt{(jq, jg)} at \((b_i,b_{*,i},\mu_i)\). The optional default
    returns \texttt{(1.0, 1.0)}.\\
order
  & \(0,1,2\)
  & Highest splitting order: LO, NLO, or NNLO, with all lower orders included.\\
method
  & \texttt{:euler} or \texttt{:rk2}
  & Explicit time integrator. The public default is \texttt{:rk2}.\\
alpha\_s\_Z\_value
  & positive Float64
  & Fixed-\(n_f=5\) reference coupling at \(91.2\,\mathrm{GeV}\); default
    \(0.118\). No threshold matching is applied.\\
alpha\_s\_loops
  & \(1,\ldots,4\)
  & Loop order used by the running coupling; default \(4\).\\
alpha\_s\_max
  & positive Float64
  & Safety ceiling for every sampled \(\alpha_s(\mu)\); default \(1.0\).\\
\bottomrule
\end{tabularx}
\end{table}

There are deliberately no \texttt{mu\_final}, \texttt{t\_start}, or
\texttt{max\_delta\_t} arguments. The complete scale interval and every time
step follow from \(b_{\min}\), \(b_{\max}\), \(N\), and \texttt{bstar\_func}.
This table documents the public \texttt{solve\_jet\_rg} entry point. The
lower-level \texttt{solve\_stepwise\_rg} function has a separate
\texttt{:euler} default and is mainly intended for testing and internal use.

\subsection{Output and interpolation}

The return value is a \texttt{StepwiseRGSolution}. It contains the two lattice
histories, boundary values, lattice metadata, and final active state. For
normal use, treat it as a callable interpolator:

\begin{lstlisting}
both = solution(1.0, 10.0)
jq = both.jq
jg = both.jg

jq_only = solution(1.0, 10.0; component = :q)
jg_only = solution(1.0, 10.0; component = :g)

canonical_mu = mu_start(solution, 1.0)
\end{lstlisting}

Valid queries satisfy
\begin{equation}
  b_{\min}\le b\le b_{\max},
  \qquad
  \ee^{t_1/2}\le\mu\le\ee^{t_N/2}.
  \label{eq:query-domain}
\end{equation}
Queries outside this rectangle raise a \texttt{DomainError}; the interpolator
does not extrapolate. At a stored lattice node it returns the stored value
exactly.

\subsection{Threading and computational cost}

Run Julia with multiple threads:
\begin{lstlisting}
julia --project="Numerical RG" --threads=auto your_script.jl
\end{lstlisting}
The implementation uses conservative parallelism:
\begin{itemize}[leftmargin=1.5em]
  \item kernel-table rows are built in parallel;
  \item independent active-node right-hand sides use at most four Julia tasks;
  \item each individual convolution dot product remains serial and uses
  \texttt{@simd}, avoiding nested threading overhead.
\end{itemize}

Running-coupling settings are stored in repository-level globals before a
kernel table is built. Consequently, solver calls with different coupling
settings must not run concurrently, and the solver should not be nested
inside another static threaded loop.

The kernel tables and solution histories require \(O(N^2)\) memory. Building
the kernel table costs \(O(N^2)\). The present explicit full-plane evolution
performs a growing convolution at many nodes and time slices, giving
\(O(N^3)\) leading work.

The benchmark script performs a warm-up and accepts the node count and BLAS
thread count:
\begin{lstlisting}
cd "Numerical RG"
julia --project=. --threads=20 test/benchmark_solver.jl 1000 1
\end{lstlisting}
Representative recorded timings with 20 Julia threads and one BLAS thread
were:
\begin{equation}
\begin{array}{c|rrrr}
N & 500 & 1000 & 2000 & 4000\\
\hline
\text{wall time [s]} & 0.45 & 3.5 & 32.6 & 286.8
\end{array}
\label{eq:timings}
\end{equation}
The hardware identifier and raw timing logs were not retained, so these values
are comparative development measurements rather than a reproducibility-grade
benchmark. The near-cubic growth is more important than the individual
numbers.

\subsection{Accuracy guidance and practical choices}

The default recommendation is \texttt{method = :rk2}. An exactly solvable
local-kernel diagnostic gives second-order convergence for RK2 and first-order
convergence for Euler on an affine boundary line. It does not test the
nonlocal channel coupling, a nonlinear \(b_*\) curve, or every detail of the
lower-region RK2 construction. The companion test report keeps the
quantitative validation results and explains these limits.

In production use:
\begin{itemize}[leftmargin=1.5em]
  \item choose \(b_{\max}\) so the jet function is already numerically zero at
  the endpoint, then enlarge it until observables are closure-insensitive;
  \item compare at least two values of \(N\);
  \item inspect the minimum canonical scale and ensure
  \(\alpha_s(\mu)\le\texttt{alpha\_s\_max}\);
  \item do not interpret agreement between two coarse grids as an analytic
  validation.
\end{itemize}

\subsection{Tests and reproducibility}

From the repository root, run
\begin{lstlisting}
julia --project="Numerical RG" -e "using Pkg; Pkg.instantiate()"

julia --project="Numerical RG" --threads=auto "Numerical RG/test/runtests.jl"

julia --project="Numerical RG" --threads=auto "Numerical RG/test/analytic_validation.jl"

julia --project="Numerical RG" --threads=auto "Numerical RG/test/momentum_sum_check.jl"
\end{lstlisting}
The regression suite contains 262 asserted checks. The two remaining scripts
are diagnostics: they print convergence and moment residuals but currently do
not fail when a numerical threshold is exceeded. Detailed logic, exact
parameters, and results are kept in \texttt{test\_report.tex}.
